\documentclass[10pt,a4paper]{amsart}
\usepackage{amsmath,amssymb,amsthm,graphicx}
\usepackage{verbatim}
\newenvironment{reference}{\comment}{\endcomment}
%\newenvironment{reference}{}{}
\newenvironment{slogan}{\comment}{\endcomment}
\newenvironment{history}{\comment}{\endcomment}

% For commutative diagrams we use Xy-pic
\usepackage[all]{xy}

% We use 2cell for 2-commutative diagrams.
\xyoption{2cell}
% Theorem environments.
%
\theoremstyle{plain}
\newtheorem{theorem}[subsection]{Theorem}
\newtheorem{proposition}[subsection]{Proposition}
\newtheorem{lemma}[subsection]{Lemma}

\theoremstyle{definition}
\newtheorem{definition}[subsection]{Definition}
\newtheorem{example}[subsection]{Example}
\newtheorem{exercise}[subsection]{Exercise}
\newtheorem{situation}[subsection]{Situation}

\theoremstyle{remark}
\newtheorem{remark}[subsection]{Remark}
\newtheorem{remarks}[subsection]{Remarks}

\numberwithin{equation}{subsection}

% Macros
%
\def\lim{\mathop{\mathrm{lim}}\nolimits}
\def\colim{\mathop{\mathrm{colim}}\nolimits}
\def\Spec{\mathop{\mathrm{Spec}}}
\def\Hom{\mathop{\mathrm{Hom}}\nolimits}
\def\Ext{\mathop{\mathrm{Ext}}\nolimits}
\def\SheafHom{\mathop{\mathcal{H}\!\mathit{om}}\nolimits}
\def\SheafExt{\mathop{\mathcal{E}\!\mathit{xt}}\nolimits}
\def\Sch{\mathit{Sch}}
\def\Mor{\mathop{\mathrm{Mor}}\nolimits}
\def\Ob{\mathop{\mathrm{Ob}}\nolimits}
\def\Sh{\mathop{\mathit{Sh}}\nolimits}
\def\NL{\mathop{N\!L}\nolimits}
\def\CH{\mathop{\mathrm{CH}}\nolimits}
\def\proetale{{pro\text{-}\acute{e}tale}}
\def\etale{{\acute{e}tale}}
\def\QCoh{\mathit{QCoh}}
\def\Ker{\mathop{\mathrm{Ker}}}
\def\Im{\mathop{\mathrm{Im}}}
\def\Coker{\mathop{\mathrm{Coker}}}
\def\Coim{\mathop{\mathrm{Coim}}}

% Boxtimes
%
\DeclareMathSymbol{\boxtimes}{\mathbin}{AMSa}{"02}

%
% Macros for moduli stacks/spaces
%
\def\QCohstack{\mathcal{QC}\!\mathit{oh}}
\def\Cohstack{\mathcal{C}\!\mathit{oh}}
\def\Spacesstack{\mathcal{S}\!\mathit{paces}}
\def\Quotfunctor{\mathrm{Quot}}
\def\Hilbfunctor{\mathrm{Hilb}}
\def\Curvesstack{\mathcal{C}\!\mathit{urves}}
\def\Polarizedstack{\mathcal{P}\!\mathit{olarized}}
\def\Complexesstack{\mathcal{C}\!\mathit{omplexes}}
% \Pic is the operator that assigns to X its picard group, usage \Pic(X)
% \Picardstack_{X/B} denotes the Picard stack of X over B
% \Picardfunctor_{X/B} denotes the Picard functor of X over B
\def\Pic{\mathop{\mathrm{Pic}}\nolimits}
\def\Picardstack{\mathcal{P}\!\mathit{ic}}
\def\Picardfunctor{\mathrm{Pic}}
\def\Deformationcategory{\mathcal{D}\!\mathit{ef}}
\usepackage{cleveref}
\crefname{originalreference}{original reference}{original references}
\Crefname{originalreference}{Original reference}{Original references}
\makeatletter
\newlabel{lemma-push-unit}{{0FHG}{source}{}{sourcepointer-1}{}}
\newlabel{lemma-push-unit@cref}{{[originalreference][0][]0FHG}{[1][0][]source}{}{}{}}
\makeatother
\crefname{originalreference}{original reference}{original references}
\Crefname{originalreference}{Original reference}{Original references}
\makeatletter
\newlabel{lemma-pushforward}{{03M9}{source}{}{sourcepointer-2}{}}
\newlabel{lemma-pushforward@cref}{{[originalreference][0][]03M9}{[1][0][]source}{}{}{}}
\makeatother
\crefname{originalreference}{original reference}{original references}
\Crefname{originalreference}{Original reference}{Original references}
\makeatletter
\newlabel{lemma-unit}{{0FHF}{source}{}{sourcepointer-3}{}}
\newlabel{lemma-unit@cref}{{[originalreference][0][]0FHF}{[1][0][]source}{}{}{}}
\makeatother
\crefname{originallemma}{original lemma}{original lemmas}
\Crefname{originallemma}{Original lemma}{Original lemmas}
\makeatletter
\newlabel{varieties-lemma-smooth-separable-closed-points-dense}{{056U}{source}{}{sourcepointer-4}{}}
\newlabel{varieties-lemma-smooth-separable-closed-points-dense@cref}{{[originallemma][0][]056U}{[1][0][]source}{}{}{}}
\makeatother
\title{Cohomology of a disjoint union of smooth projective schemes}
\author{The Stacks Project Authors}
\date{stacks-project:weil.tex:local-snapshot-20261012; GNU FDL 1.2 or later}
\begin{document}
\maketitle
\noindent\textit{Editorial note.} All delivered mathematical text is the author's own, copied from the pinned Stacks Project chapter file weil.tex; no statement or proof step was rewritten, reordered or removed.\par\medskip
\noindent\textit{Editorial note.} The counted proof is the author's complete argument: the reduction to equidimensional schemes of the same dimension, the orthogonal idempotents giving injectivity, and the base change to a finite separable extension yielding surjectivity.\par\medskip
\noindent\textit{Editorial note.} The label lemma-weil-additive names the counted statement, and the proof environment immediately following it in the source is delivered verbatim as its proof.\par\medskip
\noindent\textit{Editorial note.} The counted statement is carried by the author's own lemma environment in the chapter source.\par\medskip
\noindent\textit{Editorial note.} Pre-open gate: python operations/agent-authoring/tools/precard\_gate.py --source weil.tex --label lemma-weil-additive --start 2931 --end 3007 --local-id 28126 returned RESULT=OK (L28126 free, occupancy FREE, public interval clear). Pre-assembly tag self-check: canonical tag resolved and not present in reference/public-items.jsonl.\par\medskip
\noindent\textit{Editorial note.} The statement and proof are copied from the pinned chapter file; the cited inputs are kept as external pointers with their canonical Stacks tags.\par\medskip
\noindent\textit{Allowed standard prerequisites (editorial).}\begin{itemize}
\item Vanishing of the cohomology of the empty scheme and nonvanishing for nonempty schemes. Same chapter, the lemma computing the cohomology of the empty scheme and the unit of the theory. Handles the empty cases and provides nonvanishing of the cohomology of nonempty smooth projective schemes.
\item Pushforward along a closed and open immersion and the projection formula. Same chapter, the lemma on the pushforward of the unit and the projection formula. Shows that the classes of the two components act as orthogonal idempotents whose images give the injectivity of the comparison map.
\item Base change to a finite separable extension providing a rational point. Axiom (B) on base change for finite separable field extensions, and Varieties, the lemma that the closed points with a given residue field are dense for a smooth scheme. Produces a morphism from the second scheme to the first after base change, which yields surjectivity after tensoring.
\end{itemize}\medskip
\ifdefined\hypertarget\hypertarget{sourcepointer-1}{}\fi
\noindent\textit{Source reference.} Original reference 0FHG: sources/stacks-weil/pinned-original-source.tex, source lines 2848--2848. Prerequisite pointer invoked at proof line(s) 2972; it names a standard fact of the same chapter that the counted argument uses, and every step of the counted proof built on it is retained verbatim in the delivered body.\par\medskip
\ifdefined\hypertarget\hypertarget{sourcepointer-2}{}\fi
\noindent\textit{Source reference.} Original reference 03M9: sources/stacks-weil/pinned-original-source.tex, source lines 2660--2660. Prerequisite pointer invoked at proof line(s) 2978; it names a standard fact of the same chapter that the counted argument uses, and every step of the counted proof built on it is retained verbatim in the delivered body.\par\medskip
\ifdefined\hypertarget\hypertarget{sourcepointer-3}{}\fi
\noindent\textit{Source reference.} Original reference 0FHF: sources/stacks-weil/pinned-original-source.tex, source lines 2800--2800. Prerequisite pointer invoked at proof line(s) 2942; it names a standard fact of the same chapter that the counted argument uses, and every step of the counted proof built on it is retained verbatim in the delivered body.\par\medskip
\ifdefined\hypertarget\hypertarget{sourcepointer-4}{}\fi
\noindent\textit{Source reference.} Original lemma 056U: sources/stacks-chow/tags.txt, source lines 6722--6722. Prerequisite pointer invoked at proof line(s) 3004; it names a standard fact of a sibling chapter that the counted argument uses, and every step of the counted proof built on it is retained verbatim in the delivered body.\par\medskip
\begin{lemma}
\label{lemma-weil-additive}
Assume given (D0), (D1), (D2), and (D3) satisfying (A), (B), and (C).
If $X$ and $Y$ are smooth projective schemes over $k$, then
$H^*(X \amalg Y) \to H^*(X) \times H^*(Y)$,
$a \mapsto (i^*a, j^*a)$ is an isomorphism where $i$, $j$
are the coprojections.
\end{lemma}

\begin{proof}
If $X$ or $Y$ is empty, then this is true because
$H^*(\emptyset) = 0$ by Lemma \ref{lemma-unit}.
Thus we may assume both $X$ and $Y$ are nonempty.

\medskip\noindent
We first show that the map is injective. First, observe that
we can find morphisms $X' \to X$ and $Y' \to Y$
of smooth projective schemes so that $X'$ and $Y'$ are
equidimensional of the same dimension and such that
$X' \to X$ and $Y' \to Y$ each have a section. Namely,
decompose $X = \coprod X_d$ and $Y = \coprod Y_e$
into open and closed subschemes equidimensional of
dimension $d$ and $e$. Then take
$X' = \coprod X_d \times \mathbf{P}^{n - d}$
and $Y' = \coprod Y_e \times \mathbf{P}^{n - e}$ for some
$n$ sufficiently large. Thus pullback by
$X' \amalg Y' \to X \amalg Y$ is injective
(because there is a section) and
it suffices to show the injectivity for $X', Y'$
as we do in the next parapgrah.

\medskip\noindent
Let us show the map is injective when $X$ and $Y$ are equidimensional
of the same dimension $d$.
Observe that $[X \amalg Y] = [X] + [Y]$ in $\CH^0(X \amalg Y)$
and that $[X]$ and $[Y]$ are orthogonal idempotents in $\CH^0(X \amalg Y)$.
Thus
$$
1 = \gamma([X \amalg Y] = \gamma([X]) + \gamma([Y]) = i_*1 + j_*1
$$
is a decomposition into orthogonal idempotents. Here we have used
Lemmas \ref{lemma-unit} and \ref{lemma-push-unit} and axiom (C)(c).
Then we see that
$$
a = a \cup 1 = a \cup i_*1 + a \cup j_*1 =
i_*(i^*a) + j_*(j^*a)
$$
by the projection formula (Lemma \ref{lemma-pushforward}) and hence the map
is injective.

\medskip\noindent
We show the map is surjective. Write $e = \gamma([X])$ and $f = \gamma([Y])$
viewed as elements in $H^0(X \amalg Y)$. We have
$i^*e = 1$, $i^*f = 0$, $j^*e = 0$, and $j^*f = 1$ by axiom (C)(a).
Hence if $i^* : H^*(X \amalg Y) \to  H^*(X)$
and $j^* : H^*(X \amalg Y) \to H^*(Y)$ are surjective, then
so is $(i^*, j^*)$. Namely, for $a, a' \in H^*(X \amalg Y)$
we have
$$
(i^*a, j^*a') = (i^*(a \cup e + a' \cup f), j^*(a \cup e + a' \cup f))
$$
By symmetry it suffices to show $i^* : H^*(X \amalg Y) \to  H^*(X)$
is surjective. If there is a morphism $Y \to X$, then there is a morphism
$g : X \amalg Y \to X$ with $g \circ i = \text{id}_X$ and we conclude.
To finish the proof, observe that in order to prove
$i^*$ is surjective, it suffices to do so after tensoring
by a nonzero graded $F$-vector space. Hence by axiom (B)(b)
and nonvanishing of cohomology (Lemma \ref{lemma-unit})
it suffices to prove $i^*$ is surjective after replacing
$X$ and $Y$ by $X \times \Spec(k')$ and $Y \times \Spec(k')$
for some finite separable extension $k'/k$.
If we choose $k'$ such that there exists a closed point
$x \in X$ with $\kappa(x) = k'$ (and this is possible by
Varieties, Lemma \ref{varieties-lemma-smooth-separable-closed-points-dense})
then there is a morphism $Y \times \Spec(k') \to X \times \Spec(k')$
and we find that the proof is complete.
\end{proof}

\end{document}
